\section{稳定性与 Lyapunov 函数}
\label{sec:15-Differential-Equations-and-Dynamical-Systems/02-Stability-and-Dynamical-Systems/02}

评审会上有人问：你怎么知道这个控制器接上真实设备不会越摆越大？你手里没有闭环解，方程是非线性的，仿真只跑过几十组初值。答不上来的部分不是算力，是论证形式——你需要一个不依赖解的稳定性证明。

上一节的分类给了一半答案：算 Jacobi 矩阵的特征值，实部全为负就稳定。但那个结论有两处够不着的地方。特征值实部一旦有零，分类整个失声；而即便实部全负，Hartman--Grobman 只担保``足够近''，究竟是半径 \(0.001\) 的小球还是半径 \(100\) 的大球，它不负责。本节的工具两头都补：不依赖线性化，还能给出邻域之外的信息。

\subsection{把``稳定''拆成两个可检验的问题}

日常语言里``稳定''混着好几种意思——不炸掉、最终归位、波动不大。数学得把它拆开。

设 \(u^\ast\) 是自治系统 \(\dot u=F(u)\) 的不动点。称 \(u^\ast\) 稳定（Lyapunov 意义），若对任意 \(\varepsilon>0\) 存在 \(\delta>0\)，使得 \(\norm{u(0)-u^\ast}<\delta\) 蕴含

\[
\norm{u(t)-u^\ast}<\varepsilon
\qquad\text{对所有 }t\ge0.
\]

这个 \(\varepsilon\)-\(\delta\) 结构与连续性的定义同源。它说的是：想让轨迹永远不跑出任意小的圈子，只需把出发位置取得足够近。它防的是微小扰动被放大成巨大偏离，却不保证回到 \(u^\ast\)。无摩擦摆的轨迹永远画同样大的圆，半径等于初始偏离，稳定而不归位。

渐近稳定则要求在稳定之外，还存在一个出发邻域，其中的轨迹满足 \(u(t)\to u^\ast\)。有摩擦的摆满足这一条：摇摆越来越小，最后停在最低点。上一节的分类里各种组合都出现过——中心稳定而不渐近，稳定结点渐近稳定，鞍点与不稳定结点连普通稳定都算不上。拆开之后要问的就是两个独立问题：跑不跑得远，回不回得来。

\subsection{沿轨迹做微积分，不需要轨迹的公式}

怎样不解方程就判稳定？回到最朴素的画面：球搁在山沟底部，推一把，滚两下又回来。没有人解过那个球的运动方程，判断成立的理由是重力总把它拉向低处，而沟底已是最低点。

把``高度''抽象成状态空间上的标量函数。设 \(V(u)\) 在 \(u^\ast\) 附近连续可微，\(V(u^\ast)=0\)，且 \(u\ne u^\ast\) 时 \(V(u)>0\)，即 \(V\) 在 \(u^\ast\) 处取严格最小，称为正定。沿轨迹算它的变化率，链式法则直接给出

\[
\frac{d}{dt}V\bigl(u(t)\bigr)
=\nabla V(u)\cdot\dot u
=\nabla V(u)\cdot F(u).
\]

要点全在右端：只含已知的 \(V\) 与向量场 \(F\)，不含未知的 \(u(t)\)。若它恒 \(\le0\)，则 \(V\) 沿任何轨迹单调不增，轨迹永不爬坡。

\begin{center}
\begin{tikzpicture}[>=stealth]
  \draw[gray] (0,0) ellipse (1.9 and 1.35);
  \draw[gray] (0,0) ellipse (1.35 and 0.96);
  \draw[gray] (0,0) ellipse (0.8 and 0.57);
  \draw[gray] (0,0) ellipse (0.35 and 0.25);
  \node[gray] at (2.35,1.25) {\small \(V=c_1\)};
  \node[gray] at (1.75,-1.15) {\small \(V=c_2<c_1\)};
  \draw[very thick] plot[domain=0:14,samples=90,variable=\t]
    ({1.85*exp(-0.115*\t)*cos(\t r)},{1.31*exp(-0.115*\t)*sin(\t r)});
  \fill (0,0) circle (1.8pt);
  \node[below right] at (0.03,-0.05) {\small \(u^\ast\)};
  \draw[->] (1.85,0.05) -- (1.5,0.6);
  \node[right] at (1.9,0.45) {\small 起点};
  \node at (0,-2.0) {\small \(\dot V\le0\)：轨迹只能穿向内层，永不外出};
\end{tikzpicture}
\end{center}

\emph{正定的 \(V\) 把状态空间套成一层层等高面，\(\dot V\le0\) 说的就是轨迹只准往里穿。}

这样的 \(V\) 叫 Lyapunov 函数。它不必是物理能量，任何正定、且被动力学消耗或至少不增加的可微函数都算。绕过求解的关窍在图上一目了然：你没在算轨迹，你是在轨迹存在的前提下，用梯度与向量场的内积判断它每一刻是上坡还是下坡。非线性系统正是最需要这一手的地方——那里既没有解的公式，线性化又只在双曲点上说话，而这个内积不管方程线性与否都算得出来。

\subsection{定理只有两行，条件各管一头}

\begin{theorem}
（Lyapunov 直接法）设 \(u^\ast\) 为不动点。若在 \(u^\ast\) 的某邻域上存在正定的 \(V\)，使 \(\dot V(u)=\nabla V(u)\cdot F(u)\le0\)，则 \(u^\ast\) 稳定；若进一步对所有 \(u\ne u^\ast\) 有 \(\dot V(u)<0\)，则 \(u^\ast\) 渐近稳定。
\end{theorem}

证明的骨架直接对着定义走。在半径 \(\varepsilon\) 的球面上，连续的 \(V\) 取到正的最小值 \(m\)；由 \(V(u^\ast)=0\) 与连续性，存在 \(\delta\) 使 \(\norm{u-u^\ast}<\delta\) 时 \(V(u)<m\)；而 \(V\) 沿轨迹不增，轨迹上的函数值就永远小于 \(m\)，于是碰不到那张球面。稳定得证，\(\delta\) 就是 \(V=m\) 等高面内侧的某个半径。严格下降时再走一步：\(V\) 单调下降且有下界零，收敛到某个 \(\ell\ge0\)；若 \(\ell>0\)，轨迹滞留在 \(V=\ell\) 附近的紧集上，连续的 \(\dot V\) 在其上取到负的最大值 \(-\eta\)，下降速率被这个 \(-\eta\) 卡住，长期必然击穿 \(\ell\)，矛盾。故 \(\ell=0\)。

对一遍账。正定用来提供那个最小值 \(m\)，没有它 \(\varepsilon\) 球面上的下界可能是零，第一步就断了；\(\dot V\le0\) 用来锁住``轨迹只准往内穿''；紧性用在两处取极值，所以邻域必须是有界闭的；严格负定则专供第二段的反证。全程没有出现解的公式——存在唯一性保证轨迹合法，稳定性论证绕开求解，这是定性理论的标准姿态。

\subsection{例一：机械能恰好是阻尼摆的 Lyapunov 函数}

单摆（取单位使摆长与重力加速度归一）满足

\[
\ddot\theta+\sin\theta+b\dot\theta=0,
\qquad b\ge0,
\]

写成方程组是 \(\dot\theta=\omega,\ \dot\omega=-\sin\theta-b\omega\)。下垂平衡 \((0,0)\) 处线性化的特征值为 \((-b\pm\sqrt{b^2-4})/2\)，\(b=0\) 时纯虚，线性化沉默。

机械能

\[
E(\theta,\omega)=\frac12\omega^2+(1-\cos\theta)
\]

在原点为零、别处为正，正定。沿轨迹

\[
\dot E
=\omega\dot\omega+\sin\theta\,\dot\theta
=\omega(-\sin\theta-b\omega)+\omega\sin\theta
=-b\omega^2.
\]

两个 \(\omega\sin\theta\) 恰好抵消，剩下的系数 \(-b\) 来自阻尼，物理上就是摩擦力的功率。\(b=0\) 时 \(\dot E=0\)，能量守恒，轨迹沿 \(E\) 的等高线绕行，稳定而不渐近；纯虚特征值在这里拿到了能量解释——没有耗散，当然停不下来。

\(b>0\) 时 \(\dot E=-b\omega^2\le0\)，但它在整条 \(\omega=0\) 上都为零，只凭 \(\dot V\le0\) 拿不到渐近稳定。LaSalle 不变原理补上这个缺口：轨迹必收敛到 \(\dot E=0\) 集合中的最大不变集。在 \(\omega=0\) 上要留得住，还需 \(\dot\omega=-\sin\theta=0\)，原点附近只有 \(\theta=0\) 满足，最大不变集就是原点本身。于是有阻尼摆的下垂平衡渐近稳定。若还想要全局结论，可用 \(E\) 的径向无界性：把顶点那个不稳定平衡除外，几乎所有初值都归到下垂平衡。

\subsection{例二：目标函数就是梯度流的 Lyapunov 函数}

优化里的 Lyapunov 结构最整齐。梯度流

\[
\dot u=-\nabla f(u)
\]

沿轨迹满足

\[
\frac{d}{dt}f\bigl(u(t)\bigr)
=\nabla f(u)\cdot\dot u
=-\norm{\nabla f(u)}^2\le0,
\]

不在驻点就严格下降。目标函数自己充当 Lyapunov 函数，因为运动方向与最速下降方向重合。若 \(f\) 在极小点 \(u^\ast\) 附近严格凸，取 \(V(u)=f(u)-f(u^\ast)\) 即得渐近稳定。

这把优化和动力系统接到了一起。\hyperref[sec:09-Convex-Optimization/06-Optimization-Algorithms/01]{``梯度下降真正难在步长''}中的迭代

\[
u_{k+1}=u_k-\alpha\nabla f(u_k)
\]

正是梯度流的显式 Euler 离散。\(\alpha\) 足够小时离散迭代保住下降结构，\(f\) 兼作离散 Lyapunov 函数；\(\alpha\) 过大则以直线代曲线，一步跨过谷底飞到对面山坡，连续流本身没有的振荡发散由离散化引入。步长之难，从动力系统的角度看就是离散系统还保不保得住连续流的 Lyapunov 结构。下一节专门处理离散迭代，这个对照会再次出现。

\subsection{例三：没有现成能量时，Lyapunov 函数要拼}

不是每个系统都自带能量或目标。Lotka--Volterra 竞争模型

\[
\begin{aligned}
\dot x &= x(1-x-ay),\\
\dot y &= y(1-y-bx),
\end{aligned}
\]

中 \(a,b>0\) 是种间竞争强度，正平衡点 \((x^\ast,y^\ast)\) 满足 \(x^\ast+ay^\ast=1\) 与 \(y^\ast+bx^\ast=1\)。这里没有机械能可用，但可以拼一个出来。

注意 \(x-x^\ast\ln x\) 这个形式在 \(x=x^\ast\) 处取唯一最小，且在正象限上正定。照此构造

\[
V(x,y)=c_1\Bigl(x-x^\ast-x^\ast\ln\frac{x}{x^\ast}\Bigr)
+c_2\Bigl(y-y^\ast-y^\ast\ln\frac{y}{y^\ast}\Bigr),
\]

其中 \(c_1,c_2>0\) 待定。沿轨迹展开 \(\dot V\)，交叉项分别带着 \(c_1a\) 与 \(c_2b\)；调节两个系数的比例把交叉项消掉，例如取 \(c_1=1/a,\ c_2=1/b\)，剩下的部分为非正且只在平衡点取零，渐近稳定随之成立。这是这门手艺的常态：猜一个函数形式，用系数的自由度消掉碍事的项，而不是从口袋里掏出万能公式。

\subsection{稳定不等于准确}

设 \(e(t)\) 是室温与设定值之差，最简模型

\[
\dot e=-ke+d(t),
\qquad k>0,
\]

其中 \(-ke\) 是恒温器按误差成比例的调节，\(d(t)\) 汇总漏风、日照、烤箱这类扰动。

无扰动时取 \(V(e)=e^2/2\)，得 \(\dot V=-ke^2\le0\)，误差严格归零，原点渐近稳定。这是图纸上的性能。真实房间总有扰动。常值漏风 \(d\equiv d_0\) 让不动点偏移到 \(e=d_0/k\)：室温确实稳住了，只是没停在设定值上。若只知道 \(\abs{d(t)}\le\bar d\)，同样的分析给出

\[
\limsup_{t\to\infty}\abs{e(t)}\le\frac{\bar d}{k},
\]

误差被压进一个由扰动大小和调节强度共同决定的带子里，而不是被消灭。

所以在控制场景里必须声明：这条 Lyapunov 证明针对的是名义系统、带扰动系统，还是某个鲁棒不变集。把名义渐近稳定不加声明地当成``设备能消除任何偏差''，等于预设了一份没有证过的抗扰能力。

\subsection{找不到 \(V\) 不等于系统不稳定}

第一条边界是没有通用构造法。能量系统有机械能，梯度系统有目标函数，线性系统可以解 Lyapunov 方程 \(A^{\trans}P+PA=-Q\) 得到二次型 \(V=x^{\trans}Px\)。其余情形靠物理直觉、守恒量的组合、从已知例子改写，或者数值搜索。教材里的例子干净漂亮，实际研究中构造出一个 \(V\) 常常本身就是一篇论文的主要贡献。

第二条边界是定理单向。找到 \(V\) 就得到稳定；找不到则什么也断言不了——可能是你没找到，也可能是系统确实不稳定。没有``遍历所有可能的 \(V\)''这种操作，这个单向性是原理性的，不是熟练度问题。因此 Lyapunov 方法不能用来证明不稳定，那需要另一套工具。

第三条边界是结论的作用域跟着 \(V\) 走。\(V\) 只在某邻域上正定、\(\dot V\) 只在该邻域上非正，结论就只对邻域内的初值有效。要断言全局渐近稳定，还需 \(V\) 径向无界，即 \(\norm{u}\to\infty\) 时 \(V(u)\to\infty\)，这样等高面才在所有方向上封闭，轨迹无法沿着等高面滑向无穷。全局 Lyapunov 函数通常比局部的难构造得多。

把这三条一起记住，才算把这件工具用对：定理负责认证，不负责构造；一份 Lyapunov 证明的力量，恰好等于它所依附的那个 \(V\) 的作用域。下一节把连续时间换成离散迭代，同一套判据会以另一副面孔出现，而参数稍一变动，长期行为可以整个换掉。
